\section{Conclusion}
\label{sec:conclusion}

A critical role played by pilot studies in the process of scientific discovery is helping to determine which interventions deserve further study for potential adoption.
Motivated by this screening function and the tight budgets that such pilots often face, we argue that while representative samples are optimal in the limit for large pilots, when pilots are very resource-constrained it is optimal to design the study around a very targeted sample.
In a model with a finite number of unit types, we prove under mild conditions that for all pilot budgets below a threshold, the optimal pilot design samples only a single type.
We show that while this rule reduces to convenience sampling when the unknown treatment effects are the same across all types, under treatment effect heterogeneity the optimal pilot design prioritizes sampling from types with larger prior mean and that are more predictive of downstream effects.
Finally, we simulate a pilot study calibrated to real education experiments to illustrate that both the large- and small-budget pilot regimes can arise at realistic sample sizes.
Taken together, these findings add nuance to the common narrative that pilot studies should strive to enroll more representative samples; in fact, this can be counterproductive when budgets are tight.

A number of the work's limitations are focused on the assumptions around the cost and feasibility of sampling.
In reality, samples cannot usually be drawn uniformly at random from different subpopulations~\cite{Tipton14}, and costs are more complex than a linear function of sample size~\cite{TiptonSp21}.
Modeling more realistic mechanisms like fixed costs for access to different subpopulations, or the dynamic process of enrolling units into the study, could potentially extend the insights a model like this could provide.

Another promising direction for future work is to consider alternate continuation criteria.
While fully Bayesian approaches to decision-making have been studied, we believe that restricting the class of decision rules to those that are implementable and acceptable in science is necessary when studying pilot studies in that context.
Frequentist decision rules that reweight or extrapolate treatment effects from the study sample to the target population are one promising option, as are combining ordinary significance tests with measures of generalizability~\cite{Tipton14}.
